% Required packages:
% \usepackage{amsmath,amssymb}
% \usepackage{natbib}

\subsection{Preliminaries: Control Barrier and Lyapunov Functions}

We consider the ego control-affine system
\begin{equation}
    \dot{\mathbf{x}}
    =
    \mathbf{f}(\mathbf{x})
    +
    \mathbf{g}(\mathbf{x})\mathbf{u},
    \qquad
    \mathbf{x}\in\mathcal{X},\;
    \mathbf{u}\in\mathcal{U}.
    \label{eq:prelim_control_affine}
\end{equation}
Let \(\mathcal{V}=\{1,\ldots,N_v\}\) index surrounding vehicles,
\(\mathcal{P}=\{1,\ldots,N_p\}\) index pedestrians, and
\(\mathcal{M}=\{1,\ldots,N_m\}\) index static obstacles.  The full
constraint set at time \(t\) is
\begin{equation}
    \mathcal{A}(t)
    =
    \mathcal{V}\cup\mathcal{P}\cup\mathcal{M},
    \qquad
    N(t)=|\mathcal{A}(t)|.
    \label{eq:prelim_constraint_set}
\end{equation}

\paragraph{Control Lyapunov function.}
Let \(\mathbf{x}_g\) be the local goal.  The ego goal is encoded by
\begin{equation}
    V_{\mathrm{e}}(\mathbf{x})
    =
    \frac{1}{2}
    (\mathbf{x}-\mathbf{x}_g)^\top
    P
    (\mathbf{x}-\mathbf{x}_g),
    \qquad
    P=P^\top\succ0,
    \label{eq:prelim_clf}
\end{equation}
with the exponentially stabilizing CLF condition
\begin{equation}
    \inf_{\mathbf{u}\in\mathcal{U}}
    \left[
    \dot{V}_{\mathrm{e}}(\mathbf{x},\mathbf{u})
    +
    c_VV_{\mathrm{e}}(\mathbf{x})
    \right]
    \le0,
    \qquad
    \dot{V}_{\mathrm{e}}
    =
    L_{\mathbf{f}}V_{\mathrm{e}}
    +
    L_{\mathbf{g}}V_{\mathrm{e}}\mathbf{u},
    \label{eq:prelim_clf_condition}
\end{equation}
where \(c_V>0\).  Only the ego is controlled toward its goal; vehicles
\(i\in\mathcal{V}\) and pedestrians \(j\in\mathcal{P}\) enter through
predicted, time-varying safety constraints.

\paragraph{Control barrier function.}
Each constraint \(k\in\mathcal{A}(t)\) defines a safe set
\(\mathcal{C}_k(t)=\{\mathbf{x}:h_k(\mathbf{x},t)\ge0\}\).  For
vehicle \(i\), we use the orientation-aware barrier
\begin{equation}
    h_i(\mathbf{x},t)
    =
    1
    -
    \frac{\left(\boldsymbol{\Delta}_i^\top
    \mathbf{e}_{\parallel,i}\right)^2}{a_i^2}
    -
    \frac{\left(\boldsymbol{\Delta}_i^\top
    \mathbf{e}_{\perp,i}\right)^2}{b_i^2},
    \qquad
    \boldsymbol{\Delta}_i
    =
    \mathbf{p}(\mathbf{x})-\mathbf{p}_i(t).
    \label{eq:prelim_vehicle_barrier}
\end{equation}
For pedestrian \(j\) and static obstacle \(m\), we use
\begin{equation}
    h_j(\mathbf{x},t)
    =
    \left\|\mathbf{p}(\mathbf{x})-\mathbf{p}_j(t)\right\|_2^2-r_j^2,
    \qquad
    h_m(\mathbf{x})
    =
    \left\|\mathbf{p}(\mathbf{x})-\mathbf{p}_m\right\|_2^2-r_m^2.
    \label{eq:prelim_other_barriers}
\end{equation}
Safety requires every barrier to hold simultaneously:
\begin{equation}
    \mathcal{C}_{\mathrm{multi}}(t)
    =
    \bigcap_{i\in\mathcal{V}}\mathcal{C}_i(t)
    \cap
    \bigcap_{j\in\mathcal{P}}\mathcal{C}_j(t)
    \cap
    \bigcap_{m\in\mathcal{M}}\mathcal{C}_m(t).
    \label{eq:prelim_multi_safe_set}
\end{equation}
Equivalently, the CBF condition must hold for all \(N(t)\) constraints:
\begin{equation}
    L_{\mathbf{f}}h_k
    +
    L_{\mathbf{g}}h_k\mathbf{u}
    +
    \frac{\partial h_k}{\partial t}
    +
    \alpha_k(h_k)
    \ge0,
    \qquad
    \forall k\in\mathcal{A}(t).
    \label{eq:prelim_multi_cbf_condition}
\end{equation}
This gives the index-count relation: one condition per vehicle \(i\),
one per pedestrian \(j\), and one per static obstacle \(m\), for a total
of \(N_v+N_p+N_m\) conditions.  The single-agent case is the special
case \(N_v+N_p+N_m=1\).  The minimum safety margin and cumulative
violation are
\begin{equation}
    h_{\min}(\mathbf{x},t)
    =
    \min_{k\in\mathcal{A}(t)}h_k(\mathbf{x},t),
    \qquad
    \Phi_{\mathrm{safe}}(\mathbf{x},t)
    =
    \sum_{k\in\mathcal{A}(t)}
    \left[-h_k(\mathbf{x},t)\right]_+^2 .
    \label{eq:prelim_safety_aggregate}
\end{equation}
Thus, \(\Phi_{\mathrm{safe}}=0\) if and only if all \(N(t)\)
safe-set membership conditions \(h_k\ge0\) hold; forward invariance
additionally requires the \(N(t)\) CBF inequalities in
\eqref{eq:prelim_multi_cbf_condition}.

\paragraph{CLF-CBF quadratic program.}
The nominal control is corrected by
\begin{align}
    \mathbf{u}^{\star}
    =
    \operatorname*{arg\,min}_{\mathbf{u}\in\mathcal{U},\,\delta\ge0}
    \quad&
    \frac{1}{2}
    \left\|\mathbf{u}-\mathbf{u}_{\mathrm{nom}}\right\|_2^2
    +
    \lambda_\delta\delta^2
    \label{eq:prelim_clf_cbf_qp}
    \\
    \text{s.t.}\quad&
    L_{\mathbf{f}}V_{\mathrm{e}}
    +
    L_{\mathbf{g}}V_{\mathrm{e}}\mathbf{u}
    +
    c_VV_{\mathrm{e}}
    \le\delta,
    \nonumber
    \\
    &
    L_{\mathbf{f}}h_k
    +
    L_{\mathbf{g}}h_k\mathbf{u}
    +
    \frac{\partial h_k}{\partial t}
    +
    \alpha_k(h_k)
    \ge0,
    \quad
    \forall k\in\mathcal{A}(t).
    \nonumber
\end{align}
For SageFlow, the rectified-flow velocity supplies
\(\mathbf{u}_{\mathrm{nom}}\), the CLF term supplies goal-directed
guidance, and the CBF conditions are enforced by generation-time
projection:
\begin{equation}
    v_{\mathrm{safe}}
    =
    \operatorname{Proj}_{\mathcal{K}_{\mathrm{safe}}(t)}
    \left(
    v_{\theta}^{\mathrm{RF}}
    -
    \eta_{\mathrm{CLF}}
    \rho(\tau)\nabla V_{\mathrm{e}}
    \right),
    \label{eq:prelim_sageflow_projection}
\end{equation}
where \(\mathcal{K}_{\mathrm{safe}}(t)\) contains all \(N(t)\) CBF
conditions in \eqref{eq:prelim_multi_cbf_condition}.  Since the CLF
inequality is used as guidance rather than solved as a hard QP
constraint, it is referred to as CLF-inspired goal guidance.
